\documentclass[a4paper,12pt]{article}
\usepackage[T2A]{fontenc}
\usepackage[utf8]{inputenc}
\usepackage[english,russian]{babel}
\usepackage[left=1.5cm,right=1.5cm,top=1.5cm,bottom=1.5cm]{geometry}
% Самодостаточный файл: компилируется pdflatex без внешних зависимостей.
% Сгенерировано ЧАС ЕГЭ (fillword), зерно: fillword-e2e-42-fillword-1
% --- константы филворда ---
\newlength{\fillwordCell}
\setlength{\fillwordCell}{0.62cm} % автоподгон под A4: 8x8 клеток
\newcommand{\fillwordFont}{\fontsize{8}{10}\selectfont}
\newcommand{\fwempty}{\makebox[\fillwordCell][c]{}}
\newcommand{\fwcell}[1]{\fbox{\makebox[\fillwordCell][c]{\strut\fillwordFont #1}}}
\setlength{\fboxsep}{0pt}
\begin{document}

\section*{Филворд к варианту №1}

В сетке ниже спрятаны числа: они читаются по горизонтали (слева направо) и по вертикали
(сверху вниз) и могут пересекаться --- в общей клетке стоит одна и та же цифра.
Часть спрятанных чисел совпадает с ответами заданий этого варианта: решив вариант,
вы можете проверить себя, отыскав полученный ответ в сетке.
Список чисел специально не приводится.

\begin{center}
\begin{tabular}{*{8}{@{}c@{}}}
\fwcell{2} & \fwcell{2} & \fwcell{,} & \fwcell{0} & \fwcell{-} & \fwcell{0} & \fwcell{,} & \fwcell{9} \\
\fwcell{1} & \fwcell{4} & \fwcell{3} & \fwcell{3} & \fwcell{7} & \fwcell{,} & \fwcell{5} & \fwcell{8} \\
\fwcell{1} & \fwcell{3} & \fwcell{,} & \fwcell{5} & \fwcell{7} & \fwcell{1} & \fwcell{,} & \fwcell{4} \\
\fwcell{,} & \fwcell{9} & \fwcell{1} & \fwcell{4} & \fwcell{4} & \fwcell{6} & \fwcell{5} & \fwcell{1} \\
\fwcell{7} & \fwcell{0} & \fwcell{4} & \fwcell{0} & \fwcell{5} & \fwcell{7} & \fwcell{,} & \fwcell{4} \\
\fwcell{7} & \fwcell{-} & \fwcell{4} & \fwcell{-} & \fwcell{5} & \fwcell{7} & \fwcell{5} & \fwcell{3} \\
\fwcell{,} & \fwcell{0} & \fwcell{1} & \fwcell{7} & \fwcell{,} & \fwcell{4} & \fwcell{-} & \fwcell{1} \\
\fwcell{8} & \fwcell{6} & \fwcell{0} & \fwcell{5} & \fwcell{4} & \fwcell{2} & \fwcell{2} & \fwcell{5}
\end{tabular}
\end{center}

\newpage

\section*{Филворд к варианту №1: ключ для учителя}

Расположение слов (ключ не предназначен для печати в варианте ученика):

\begin{itemize}
\item 13 --- по горизонтали, начало в клетке (1; 3);
\item 8 --- по горизонтали, начало в клетке (1; 8);
\item 3 --- по горизонтали, начало в клетке (8; 6);
\item 337,5 --- по горизонтали, начало в клетке (3; 2);
\item 98 --- по вертикали, начало в клетке (8; 1);
\item 0,167 --- по вертикали, начало в клетке (6; 1);
\item 2 --- по горизонтали, начало в клетке (1; 1);
\item 3,5 --- по горизонтали, начало в клетке (2; 3);
\item -0,9 --- по горизонтали, начало в клетке (5; 1);
\item 6 --- по горизонтали, начало в клетке (2; 8);
\item 225 --- по горизонтали, начало в клетке (6; 8);
\item 15 --- по вертикали, начало в клетке (8; 7);
\item 243 --- по вертикали, начало в клетке (2; 1);
\item 390 --- по вертикали, начало в клетке (2; 3);
\item -1 --- по горизонтали, начало в клетке (7; 7);
\end{itemize}

В сетке также присутствуют числа-обманки (2 шт.): они стоят на поле, но к ответам варианта не относятся и в ключе не перечислены.

\end{document}
